\subsubsection{Computational Time Requirements}
\label{sec:comp_req}

To maintain audio quality in real-time, the deadline must meet the processing requirements, the audio engine
needs to finish processing to synthesize audio and apply DSP within the buffer period \cite{zhao2024multithreaded}, 
until the next chunk of samples are processed. Derived from Rate Monotonic Scheduling theory (\Cref{prompt:execution_time}) 
\cite{openai_chatgpt}, the ratio ($U$) between the callback execution time, $C$ and the buffer duration ($t_{\text{buffer}}$)
determines the the real-time processor utilization, representing the portion of the available deadline 
consumed by the audio engine. This is described by the following equation below \cite{liu1973scheduling}:

\begin{equation}\label{eq:execution_time}
    U = \frac{C}{t_{\text{buffer}}}
\end{equation}

The processing time satisfies the deadline requirements, if and only if $C < t_{\text{buffer}}$ or $U < 1$.
Otherwise, it will lead to glitches (observed through clicks and pops) and some unwanted artifacts.

Design a real-time friendly audio engine with DSP, it must maintain real-time constraints, which includes
avoiding any function calls that takes a longer period of time to execute. These might include the following \cite{portaudio}:
\begin{itemize}
    \item Memory allocation/deallocation
    \item I/O (e.g. \texttt{printf()}, \texttt{std::cin}, etc.)
    \item Context switching (e.g. \texttt{exec()}, or \texttt{yield()})
    \item Any mutex operations
    \item Any tasks related to the operating system
    \item Functions with too much time complexity
\end{itemize}

\Cref{app:exec_time} shows the code listings which measures the callback execution times using \texttt{std::chrono::high\_resolution\_clock::now()}. 
It takes multiple readings for as long as the audio engine runs, which are useful for statistical analysis. This includes the mean, 
maximum (worst-case execution time), and the standard deviation (determines consistency). If the mean or maximum callback time is greater 
than the buffer duration, then there are glitches present during the stream, so it either is a performance issue or a bug needs to 
be fixed.

A numerically different approach for calculating the standard deviation is conducted using the accumulation 
of the sum of callback times, and the squared callback times on to a single variable,
instead of an array or \texttt{std::vector}. The following equation determines this standard deviation \Cref{prompt:stddev}
\cite{haas2022sumofsquares,openai_chatgpt}:

\begin{equation}\label{eq:stddev}
    \sigma = \sqrt{ \frac{\sum x^2 - \frac{\left(\sum x\right)^2}{n} }{n - 1} }
\end{equation}

A larger standard deviation indicates larger variations in time, which means more jitter.

\Cref{app:exec_time}, Listing \ref{lst:raii-timer} shows another way of measuring time, which utilizes the RAII (Resource Acquisition Is Initialization)
technique by taking the start and end times at the constructor and destructor, and calculating the elasped time
after the object is destroyed \cite{nagi2026timingcpp}.
